\chapter{Information-theoretic tools}
\label{chap:prereq_information}

These lemmas are in \texttt{LMLPapers} (\texttt{ForMathlib/InformationTheory/KullbackLeibler/},
\texttt{Identification/ChangeOfMeasure.lean}) or in LML
(\texttt{SequentialLearning/DivergenceDecomposition.lean}), and are to be copied or adapted.

\begin{lemma}[Bretagnolle--Huber inequality]
  \label{lem:bretagnolle_huber}
  \lean{InformationTheory.bretagnolle_huber}
  \leanok
  For probability measures $\mu, \nu$ and a measurable set $E$,
  $\mu(E) + \nu(E^c) \ge \frac12 \exp(-\KL(\mu \,\|\, \nu))$.
\end{lemma}

\begin{proof}
  \leanok
  $\mu(E) + \nu(E^c) \ge \int \min(d\mu, d\nu) \ge \frac12 (\int \sqrt{d\mu\, d\nu})^2 \ge \frac12
  \exp(-\KL(\mu \,\|\, \nu))$, by the Cauchy--Schwarz inequality and Jensen's inequality.
\end{proof}

\begin{lemma}[Divergence between Gaussian laws]
  \label{lem:kl_gaussian}
  \lean{InformationTheory.klDiv_gaussianReal_one, InformationTheory.klDiv_gaussianReal_same_var,
    InformationTheory.klDiv_gaussianReal}
  \leanok
  $\KL(\Gauss(m, 1) \,\|\, \Gauss(m', 1)) = (m - m')^2 / 2$. In particular the reward laws of the
  arm $x$ in the linear bandits of parameters $\theta$ and $\theta'$ with standard Gaussian noise
  have divergence $(x^\top(\theta - \theta'))^2 / 2$.
\end{lemma}

\begin{proof}
  \leanok
  Direct computation of the integral of the log-likelihood ratio, an affine function.
\end{proof}

\begin{lemma}[Data processing for the output]
  \label{lem:data_processing_output}
  \lean{Learning.IdentAlg.IsRun.klDiv_map_out_le,
    Learning.IdentAlg.IsFixedBudget.klDiv_outputMeasure_le}
  \leanok
  For an identification algorithm with budget $T$ and runs against two environments, the divergence
  between the laws of the outputs is at most the divergence between the laws of the histories of
  the first $T$ rounds.
\end{lemma}

\begin{proof}
  \leanok
  The output is drawn from the same kernel applied to the history of the $T$ rounds; data
  processing inequality for kernels.
\end{proof}

\begin{lemma}[Chain rule for non-stationary bandits]
  \label{lem:chain_rule_nonstationary}
  \lean{Learning.IsAlgEnvSeq.klDiv_map_history_banditSeq_eq_sum,
    Learning.IsAlgEnvSeq.klDiv_map_history_banditSeq,
    Learning.IsAlgEnvSeq.klDiv_map_history_compProd_banditSeq, Learning.stepKernel_banditSeq}
  \leanok
  For runs of one algorithm against the non-stationary bandits $\nu$ and $\nu'$, the divergence
  between the laws of the histories of the first $T$ rounds is
  $\sum_{t < T} \sum_a \Prob_\nu(X_t = a)\, \KL(\nu_{t, a} \,\|\, \nu'_{t, a})$ (finitely many arms).
\end{lemma}

\begin{proof}
  \leanok
  LML's chain rule for histories (\texttt{IsAlgEnvSeq.klDiv\_map\_history\_stepKernel}): the
  divergence of the step of round $t$ given the history is that of the reward kernels, the policy
  kernels being equal; the divergence of $\Prob_\nu(X_t \in \cdot) \otimes \nu_{t, \cdot}$ and
  $\Prob_\nu(X_t \in \cdot) \otimes \nu'_{t, \cdot}$ is the average of the divergences.
\end{proof}
